%%-----------------------------%%
%%        STAKEHOLDERS         %%
%%-----------------------------%%

\section{Stakeholders}

\begin{explicacao}[Importante]
Um mesmo usuário pode possuir múltiplos papéis, observadas as combinações permitidas na Matriz de Papéis e Permissões da seção de Regras de Negócio. O papel \textit{Chefe Geral} é mutuamente exclusivo com os demais papéis para uma mesma conta. O papel \textit{Bolsista} possui como pré-condição existencial e estrutural o papel de \textit{Aluno} ($Bolsista \implies Aluno$).
\end{explicacao}

%%--------------------------------------------%%
%%       STAKEHOLDERS - CHEFE GERAL           %%
%%--------------------------------------------%%
\subsection{Chefe Geral}
Esse é o maior cargo do sistema: Esse cargo gerencia os outros cargos.
\begin{itemize}
    \item Poder cadastrar outros usuários e definir os papéis deles no sistema, de acordo com as regras descritas em \textit{Regras de Negócio} e na \textit{Descrição das telas}.
    \item Pode cadastrar um Almoxarifado. 
    \item Ao cadastrar um Gestor de Almoxarifado, pode colocá-lo responsável por mais de um almoxarifado. 
    \item Ao cadastrar um Gestor de Almoxarifado, pode vinculá-lo imediatamente para gerenciar um ou mais almoxarifados, ou apenas cadastrar o usuário (deixando a vinculação para um momento futuro, resultando em uma dashboard sem almoxarifados ativos inicialmente).
    \item Atribuir o papel de Bolsista a um usuário que já possua o papel Aluno.
\end{itemize}

%%-----------------------------------------------%%
%%     STAKEHOLDERS - GESTOR DE ALMOXARIFADO     %%
%%-----------------------------------------------%%
\subsection{Gestor de Almoxarifado}
Pode haver mais de um gestor, ele(s) pode(m):
\begin{itemize}
    \item Gerar um pdf de etiquetas de frasco 
    \item Adicionar/alterar a quantidade de um reagente no estoque.
    \item Pesquisar por um reagente no estoque.
    \item Cadastrar uma nova especificação de reagente, informando os dados obrigatórios definidos pelo tipo de substância e pelo estado físico.
    \item Marcar um reagente como removido (não exclui do banco, apenas muda o status).
    \item Alterar o status de um reagente: marcar empréstimo, devolução (através de pesagem), descarte, vazio ou quebrado.
    \item Autorizar expressamente o uso de frascos vencidos para fins didáticos, demonstrações e pesquisa excepcional conforme Q04.
    \item Gerar relatório mensal (ou de um período específico) para um almoxarifado ou para todos.
    \item Editar/adicionar propriedades de um reagente.
\end{itemize}

%%-----------------------------------------------%%
%%  STAKEHOLDERS - GESTOR DE BENS PATRIMONIAIS   %%
%%-----------------------------------------------%%
\subsection{Gestor de Bens Patrimoniais}
Pode haver mais de um gestor de bem patrimonial, ele(s) pode(m):
\begin{itemize}
    \item Visualizar as requisições dos professores para adicionar/editar um bem patrimonial.
    \item Quando a requisição do professor é para colocar um bem como inservível e o gestor aprova, o gestor deve dar baixa do bem patrimonial através do processo da UENF (abrir processo no SEI etc.). Após esse processo, ele marca no sistema que foi dada baixa.
    \item Adicionar um bem patrimonial e criar resumos de catálogo.
    \item Editar um bem patrimonial e suas propriedades.
    \item Gerar e baixar relatório de bens patrimoniais mensal (ou de período específico).
    \item Visualizar os bens patrimoniais existentes.
    \item Pesquisar bens patrimoniais por número de patrimônio ou nome.
    \item Filtrar bens patrimoniais por local (prédio, andar e sala), estado de conservação ou status.
    \item Imprimir Relatório mensal (ou período personalizado), ou sobre um Bem patrimonial específico podendo filtrar por prédio, andar, sala, Nome do responsável, etc.
\end{itemize}

%%-----------------------------------%%
%%        STAKEHOLDERS - PROFESSOR   %%
%%-----------------------------------%%
\subsection{Professor}
Com certeza haverá mais de um, eles podem:
\begin{itemize}
    \item Criar uma turma (especificando nome e capacidade máxima de alunos).
    \item Copiar o código de uma turma dele.
    \item Criar um post em uma turma (ao criar um post, o professor pode selecionar um roteiro que ele já tenha feito o upload, um roteiro que outro professor compartilhou com ele, ou ainda fazer upload de um roteiro do próprio computador -- isso abre um modal para cadastrar o roteiro).
    \item Adicionar usuários apenas como alunos por e-mail (mandando o link para eles trocarem a senha).
    \item Pesquisar por bens patrimoniais.
    \item Abrir uma requisição de adição ou edição de bem patrimonial (analisada por um gestor de bens patrimoniais).
    \item Adicionar novos Locais (Prédio, Andar e Sala) no sistema de forma autônoma (por exemplo, ao fazer uma requisição, selecionar a opção "Novo Local" e registrá-lo).
    \item Solicitar/receber empréstimo de reagentes do almoxarifado; o registro operacional da retirada é efetuado pelo Gestor de Almoxarifado.
    \item Adicionar um roteiro de experimento (formato PDF).
    \item Vincular um roteiro de experimento a uma turma.
    \item Compartilhar um roteiro de experimento com outros professores.
    \item Acessar a aba de roteiros compartilhados (dividida em: ``Compartilhados comigo'' e ``Eu compartilhei'').
    \item Baixar um roteiro (seja adicionado por ele mesmo ou compartilhado com ele por outro professor).
    \item Vincular um roteiro de experimento compartilhado com ele por outro professor a uma turma dele.
    \item Arquivar e desarquivar turmas sob sua responsabilidade.
\end{itemize}

%%-----------------------------------%%
%%        STAKEHOLDERS - ALUNO       %%
%%-----------------------------------%%
\subsection{Aluno}
Esses com certeza serão maioria. Eles podem:
\begin{itemize}
    \item Entrar em uma turma usando o código disponibilizado pelo professor (respeitado o limite de capacidade da turma).
    \item Visualizar suas turmas e o professor respectivo de cada turma.
    \item Baixar qualquer roteiro de qualquer turma da qual façam parte.
    \item Visualizar os colegas que estão na mesma turma.
    \item Comentar em um post (o comentário é visível tanto para o professor quanto para os colegas).
\end{itemize}

\subsection{Bolsista}
Seja de graduação ou pós-graduação, este papel identifica um usuário que possui o papel de Aluno e uma autorização adicional para retirada de reagentes ($Bolsista \implies Aluno$). Ele é mutuamente compatível com Aluno, mas possui restrições próprias para acumulação com papéis de gestor.
\begin{itemize}
    \item Ele pode ver os reagentes e seus estoques assim como os professores.
    \item O gestor de reagentes pode registrar retiradas e devoluções para esses usuários.
    \item Bolsista e Gestor de Almoxarifado são mutuamente exclusivos (SoD). Aluno sem Bolsista pode exercer Gestor de Almoxarifado, conforme DP-C01.
\end{itemize}

\subsection{Matriz de Papéis e Permissões}
Esta tabela é um resumo de stakeholders e UX. A fonte normativa única da
autorização, inclusive estado ativo, escopo, ownership, decisão transacional,
revogação e resultado de negação, é a matriz M9 da Seção~7. Papéis exibidos ou
informados pelo cliente nunca determinam permissão.


\small
\setlength{\tabcolsep}{4pt}
\small
\setlength{\tabcolsep}{3pt}
\begin{longtable}{>{\raggedright\arraybackslash}p{0.24\textwidth}*{6}{>{\centering\arraybackslash}p{0.105\textwidth}}}
\toprule
\textbf{Ação} & \textbf{Chefe} & \textbf{Gestor Almox.} & \textbf{Gestor Patr.} & \textbf{Professor} & \textbf{Aluno} & \textbf{Bolsista}\\
\midrule
\endfirsthead
\toprule
\textbf{Ação} & \textbf{Chefe} & \textbf{Gestor Almox.} & \textbf{Gestor Patr.} & \textbf{Professor} & \textbf{Aluno} & \textbf{Bolsista} (cont.)\\
\midrule
\endhead
\bottomrule
\endfoot
Gerenciar usuários e papéis & \checkmark & -- & -- & -- & -- & --\\
Gerenciar almoxarifados & \checkmark & -- & -- & -- & -- & --\\
Consultar estoque de reagentes & \checkmark & \checkmark & -- & \checkmark & \checkmark & \checkmark\\
Cadastrar/editar reagente e frasco & \checkmark & \checkmark & -- & -- & -- & --\\
Registrar retirada/devolução (operação de balcão) & \checkmark & \checkmark & -- & retirante & -- & retirante\\
Gerenciar bens patrimoniais & \checkmark & -- & \checkmark & leitura & -- & --\\
Abrir requisição de patrimônio & \checkmark & -- & -- & \checkmark & -- & --\\
Responder requisição de patrimônio & \checkmark & -- & \checkmark & -- & -- & --\\
Gerenciar turmas próprias & \checkmark & -- & -- & \checkmark & -- & --\\
Ingressar em turma & -- & -- & -- & -- & \checkmark & \checkmark\\
Criar/comentar posts permitidos & -- & -- & -- & \checkmark & \checkmark & \checkmark\\
Gerenciar roteiros próprios & \checkmark & -- & -- & \checkmark & -- & --\\
Compartilhar roteiros & \checkmark & -- & -- & \checkmark & -- & --\\
Cadastrar/editar matérias & \checkmark & -- & -- & \checkmark & -- & --\\
Gerar etiquetas de frasco & \checkmark & \checkmark & -- & -- & -- & --\\
\bottomrule
\end{longtable}


\begin{explicacao}[Regra de autorização]
A Cloud Function responsável por cada mutação deve resolver os papéis a partir dos registros persistidos e, quando aplicável, validar também o escopo do recurso. Por exemplo, um Gestor de Almoxarifado só atua sobre almoxarifados aos quais está vinculado. O frontend pode ocultar ou desabilitar ações, mas não constitui autorização.

A linha ``Gerenciar turmas próprias'' pertence ao professor dono. O Chefe Geral não possui turmas próprias (RN-ROLE-01) e não recebe ownership de turma alheia: sua atuação acadêmica é a intervenção excepcional de moderação Q13, com justificativa, auditoria e a marcação institucional nas publicações. O contrato de Turma, matrícula e convites é a Seção~\ref{sec:regras-turmas-m11}.
\end{explicacao}

\paragraph{Operador e retirante (AUDITORIA-7, C002).}
A autorização de Chefe Geral na matriz refere-se a operar o balcão, identificado por \texttt{id\_gestor\_retirada}. O destinatário \texttt{id\_usuario\_retirou} deve ser Professor ou Bolsista ativo. Chefe Geral não pode ser retirante, pois seu papel é exclusivo (RN-ROLE-01).
